\section{Signed weights and the limits of mean criteria}\label{sec:signed}
Write $R_w=G^*W G$ and $Q_w(A)=\ip{A}{K_{0,n}A}+\ip{A}{R_wA}$.
For signed weights this form need not be nonnegative. Its spatial means
alone no longer supply the positive trace comparison used above.
The following sufficient condition is nevertheless available without a
lower bound on any original weight. Define the \emph{pointwise} positive part
and its selected directional mean by
\[
 w^+(k,(x,l))=\max\{w(k,(x,l)),0\},\qquad
 b^+_{ij}(n)=n^{-d}\sum_x w^+(i,(x,j)).
\]
In general $b^+_{ij}$ is not the positive part of the signed mean.

\begin{theorem}[A sufficient signed-weight condition]\label{thm:signed-positive}
Keep the observation, penalty, local divergence-free domain and fixed
$c>0$ of Theorem~\ref{thm:average}, but use arbitrary signed actual weights
in the local coercivity inequality. For every $0\leq R<n$,
\[
 c(2R-\dim F_n)\leq
       \frac{8\pi^2(1+b^+_{ij}(n))R^3}{n^2}.
\]
If $(1+b^+_{ij}(n))/n^2\to0$, then $\dim F_n\to\infty$.
The original $C_n,L_n,\Phi_n,c,K_n,\delta_n,r_n$ are retained.
\end{theorem}
\begin{proof}
The exact weighted-gradient energy gives, for every actual cochain,
\[
 \ip{A}{R_w A}=\sum_{x,k,l}w(k,(x,l))\|D_k A_l(x)\|^2
       \leq\sum_{x,k,l}w^+(k,(x,l))\|D_k A_l(x)\|^2
       =\ip{A}{R_{w^+}A}.
\]
The inequality uses only nonnegativity of the squared differences.
Thus the original local coercivity inequality transfers upward to the
constructed nonnegative repair $R_{w^+}$, and Theorem~\ref{thm:average}
applies. Nonnegativity of $w^+$ and the comparison are proved from their
actual definitions, not supplied as fields of a comparison structure.
\end{proof}
This changes the explicit mean hypothesis for signed weights. It neither
proves the original weights nonnegative nor removes the local regularity,
quadratic penalty or fixed positive coercivity requirements.

\subsection{Zero signed means with fixed output}
There is no counterpart using only ordinary signed directional means,
even when \emph{every} such mean is zero. Fix a site $o\in X_n$, set
$V=n^d$, $a=2dn^2$, and construct actual weights
\[
 w_n(k,(x,l))=a\bigl(1-V\mathbf1_{x=o}\bigr).
\]
Since $|X_n|=V$, their ordinary spatial means are all exactly zero.
Let $\mu(A)$ be the directional mean above, and let
$J_n(A)=(D_k A_l(o))_{k,l}$ collect all actual forward differences at $o$.
Use the same finite output type for every side:
\[
 F=(\R^m)^d\times(\R^m)^{d\times d},\quad
 C_n(A)=(\mu(A),J_n(A)),\quad
 \|(y,z)\|_F=\max\{\|y\|_2,\|z\|_2\}.
\]
Its dimension is $(d+d^2)m$, independent of $n$. The companion code proves
the fixed finite output type; this displayed numerical dimension is an
analytic calculation, not an additional named oracle endpoint.

\begin{proposition}\label{prop:signed-zero}
The actual continuous linear observation $C_n$ and the nonnegative
quadratic penalty
\[
 \Phi_n(y,z)=2V\|y\|_2^2+aV\|z\|_2^2
\]
satisfy $C_n(0)=0$, $DC_n(0)=C_n$,
$\Phi_n(y,z)\leq V(2+a)\|(y,z)\|_F^2$, and
\[
 \|A\|^2\leq Q_{w_n}(A)+\Phi_n(C_n(A))\qquad(A\in E_n).
\]
Consequently all the original local observation and penalty hypotheses
can hold with fixed output and $c=1$, although all signed means vanish.
\end{proposition}
\begin{proof}
Linearity follows directly from the finite sums and differences. The
penalty bound follows from each projection norm being at most the
specified product max norm. The exact energy expansion is
\[
 \ip{A}{R_{w_n}A}=a\|GA\|^2-aV\|J_n(A)\|_2^2.
\]
Adding the penalty cancels the negative term and leaves
\[
 Q_{w_n}(A)+\Phi_n(C_n(A))
   =\ip{A}{K_{0,n}A}+a\|GA\|^2+2V\|\mu(A)\|_2^2
   \geq\|A\|^2.
\]
The last step is the actual mean/path estimate of the preceding section,
with $a=2dn^2$, together with $K_{0,n}\succeq0$.
\end{proof}
The coefficient $K_n=V(2+a)$ may grow with the lattice side, exactly as
allowed by the proposed ordinary-mean extension. This is not a transfer
of the desired coercivity into an assumed observation certificate.
The weights, observation, penalty and cancellation are constructed.
For $n>1$ their positive-part mean is $a(1-V^{-1})$, of quadratic rather
than subquadratic order, so they do not contradict
Theorem~\ref{thm:signed-positive}.

As an analytic diagnostic, take $n\geq3$, $i\ne j$, a unit internal vector
$u$, and $A_j(x)=u\mathbf1_{x_i=o_i}$, with other components zero. Then
$\delta A=0$, $\|GA\|^2=2n^{d-1}$ and $\|J_n(A)\|^2=1$, giving
\[
 Q_{w_n}(A)=n^{d-1}\bigl(2-a(n-2)\bigr)<0.
\]
This explains the failure of the original positive trace argument.
The diagnostic field calculation is analytic; the formal fixed-output
family theorem does not claim to check this additional field explicitly.

\subsection{An attributed square cancellation}
Fix distinct directions $i,k$, a component $j\ne i$ (for instance $j=k$),
and $o\in X_n$. For an amplitude $A_0\geq0$, define the repair
weights additively by
\begin{align*}
 w(p)&=A_0\mathbf1_{p=(i,(o,j))}
       -3A_0\mathbf1_{p=(k,(o,j))}\\
     &\quad-3A_0\mathbf1_{p=(i,(o+e_k,j))}
       -3A_0\mathbf1_{p=(k,(o+e_i,j))}.
\end{align*}
Writing $D_s A_j(x)=A_j(x+e_s)-A_j(x)$, commuting shifts give
\[
 D_i A_j(o)=D_k A_j(o)+D_i A_j(o+e_k)-D_k A_j(o+e_i).
\]
The squared norm of this sum is at most three times the sum of the
three squared norms. The actual finite weighted-energy sum therefore
proves $\ip{A}{R_wA}\leq0$ for every cochain.
For $n\geq2$ the positive edge is separate from all negative contributions,
so $w^+=A_0\mathbf1_{p=(i,(o,j))}$ and $b^+_{ij}=A_0/n^d$.
Choosing $A_0=n^d n^2$ gives $b^+_{ij}=n^2$.

If the original local coercivity hypotheses hold for this actual repair,
they imply the same hypotheses with zero repair, since $Q_w(A)\leq
\ip{A}{K_{0,n}A}$. Applying Theorem~\ref{thm:average} with zero weights
again forces $\dim F_n\to\infty$. The selected normalized positive-part mean is
$1+n^{-2}$ and does not tend to zero. This last limit interpretation is
an elementary analytic consequence of the formal exact mean identity,
not a separate negated-limit oracle endpoint. This single sparse cycle
does not make every directional criterion fail: other eligible means
are zero. Its formal endpoints establish the selected mean and derived
order, not an all-directions counterexample.

There is, however, a finite sum that covers every eligible pair. This
extension is analytic here. For every $j$ and every $i\ne j$, use the
preceding cycle with $k=j$, the same site $o$ and amplitude $A_0$, and add
all its weight functions to obtain $\overline w$. Linearity of the exact
weighted-energy sum gives $\ip{A}{R_{\overline w}A}\leq0$.
For every off-diagonal pair $i\ne j$, direct inspection of the four
indicators gives
\[
 \overline w(i,(x,j))=
 A_0\mathbf1_{x=o}-3A_0\mathbf1_{x=o+e_j}.
\]
Every diagonal component has only nonpositive contributions. For $n\geq2$
the two displayed sites are distinct, so each eligible pair has
$\overline b^+_{ij}=A_0/n^d$. With $A_0=n^d n^2$, \emph{all} off-diagonal
positive-part means equal $n^2$, while the derived domination by zero
still gives the original rank obstruction. Thus no eligible mean condition
in Theorem~\ref{thm:signed-positive} applies to this sum. The conclusion
follows from its actual energy order, not from an assumed comparison
field. The finite sum and its all-pairs mean calculation are not silently
included in the single-cycle formal oracle reports.

The cancellation mechanism is classical, rather than a new general
spectral principle. Restricted to the four active vertices, the negative
of this scalar repair has three positive path conductances $3A_0$ and a
direct negative conductance $-A_0$. For $A_0>0$ the path resistance is
$1/A_0$, so this is the critical case of Zelazo and B\"urger's
Theorem III.3~\cite{zelazo}. Isolated vertices contribute zero blocks;
the internal vector coordinates sum scalar copies. This comparison
attributes the mechanism; it does not establish originality of the
whole consolidated dossier.

The conditional obstruction is nonvacuous. An analytic check takes
$F_n=E_n$, $C_n=I$ and
$\Phi_n(y)=\|y\|^2-\ip{y}{R_w y}\geq0$ (or replace $w$ by $\overline w$).
The operator norm gives the upper quadratic bound
$(1+\|R_w\|)\|y\|^2$, while $Q_w(A)+\Phi_n(A)
=\ip{A}{K_{0,n}A}+\|A\|^2\geq\|A\|^2$.
This example has growing output, as the theorem requires; it is not an
additional formal endpoint or a compressed observation construction.
None of these signed-weight statements proves a volume-density bound,
identifies an interacting Hessian, or supplies a physical mass gap.
